\chapter{\hfill{\centering Entreprise d'acceuil}\hfill} %Présentation de l'entre...
\newpage
%----------------------------------------------------------------------
%           Presentation de l’entreprise BEWIZE
%----------------------------------------------------------------------
\section{Introduction}
Ce chapitre a pour objet de présenter le startup bewize. D’abord il définit l’organisme d’accueil, traite la problématique et présente les objectifs du projet. Enfin il décrit la méthodologie et le planning du projet. 


\section{Presentation de l’entreprise BEWIZE}
\quad \textbf{BEWIZE} est une start-up Edtech ( Educational technology) créée en 2021 qui digitalise le contenu pédagogique du primaire tout en assurant un suivi personnalisé par élève . Bewize une startup qui cherche à révolutionner l'éducation en proposant une approche innovante et ludique pour aider les élèves à atteindre leur plein potentiel. Grâce à son application mobile, son back-office et ses outils d'apprentissage interactifs, Bewize offre une expérience d'apprentissage unique qui favorise l'autonomie et l'épanouissement des élèves.



\begin{figure}[h]
\begin{center}
            
\includegraphics[width=0.5\textwidth]{Logos/bewize.jpg}\\[0.5cm]
\end{center}
\caption{Logo de BEWIZE}
\end{figure}

Le produit principal de Bewize est l'application mobile Bewize, qui est conçue pour aider les élèves à réviser leurs cours de manière efficace et ludique. Cette application est accompagnée d'un back-office, le tableau de bord école, qui permet aux enseignants de suivre les progrès de leurs élèves et de leur fournir un suivi personnalisé ainsi que un autre back-office pour la gestion de contenu au sein du produit principale Bewize.


%----------------------------------------------------------------------
%           Pourquoi Bewize ?
%---------------------------------------------------------------------
\section{Pourquoi Bewize ?}
\subsection{Contenu conforme au référentiel du ministère}
Les contenus proposés par \textbf{BEWIZE} sont conformes au référentiel établi par le ministère de l'Éducation. Pour garantir cela, l'entreprise a formé une équipe composée de plus de 30 enseignants et experts pédagogiques de l'Éducation Nationale. Ces professionnels travaillent en étroite collaboration pour élaborer des contenus pédagogiques de qualité, adaptés aux programmes scolaires actuels et aux dernières réformes éducatives.

Cette approche permet à \textbf{BEWIZE} de s'assurer que ses contenus sont pertinents, actualisés et alignés sur les objectifs d'apprentissage définis par le ministère de l'Éducation. Ainsi, les élèves bénéficient d'un contenu éducatif fiable et complet, garantissant une expérience d'apprentissage de haute qualité.

\subsection{Parcours conçu sur la base d'une pédagogie innovante}
\textbf{BEWIZE} a conçu un parcours d'apprentissage interactif qui combine différentes techniques pédagogiques innovantes et éprouvées. L'objectif est de créer une expérience d'apprentissage engageante et efficace pour les élèves. Voici quelques-unes des techniques utilisées :

Apprentissage ludique : \textbf{BEWIZE} intègre des éléments ludiques tels que des jeux éducatifs, des quiz interactifs et des défis pour rendre l'apprentissage plus amusant et motivant. Cela permet aux élèves d'apprendre de manière plus engageante et de renforcer leur intérêt pour les sujets étudiés.

Apprentissage adaptatif : La plateforme de \textbf{BEWIZE} utilise des algorithmes d'apprentissage automatique pour personnaliser le contenu en fonction des besoins et du niveau de chaque élève. Elle adapte le parcours d'apprentissage en proposant des activités et des ressources correspondant au niveau de compétence de chaque élève, favorisant ainsi une progression adaptée et personnalisée.

Rétroaction immédiate : \textbf{BEWIZE} offre une rétroaction immédiate aux élèves, leur permettant de comprendre leurs erreurs et de corriger leurs lacunes rapidement. Cette rétroaction peut prendre la forme de commentaires détaillés, d'explications supplémentaires ou de suggestions d'amélioration, ce qui favorise une compréhension approfondie des sujets abordés.

Collaboration et partage : La plateforme de \textbf{BEWIZE} encourage la collaboration entre les élèves en leur permettant de partager leurs idées, de travailler ensemble sur des projets et de participer à des activités en groupe. Cela favorise l'apprentissage social, le développement des compétences de communication et la création d'un environnement d'apprentissage interactif.

En combinant ces différentes techniques pédagogiques, \textbf{BEWIZE} vise à offrir aux élèves une expérience d'apprentissage complète, interactive et efficace, favorisant ainsi leur autonomie et leur épanouissement.

\subsection{Apprentissage autonome accessible à tout moment}
Avec \textbf{BEWIZE}, les parents peuvent accéder à l'expérience d'apprentissage de leur enfant directement depuis leur smartphone ou leur tablette. L'application mobile de \textbf{BEWIZE} offre une interface conviviale et intuitive qui permet aux parents de suivre les progrès de leur enfant et d'interagir avec son parcours d'apprentissage.

Grâce à cette application, les parents peuvent :

\textbf{Suivre les progrès }: Les parents ont un accès en temps réel aux résultats et aux performances de leur enfant. Ils peuvent consulter les activités terminées, les scores obtenus et les compétences acquises. Cela permet aux parents de suivre de près les progrès de leur enfant et de comprendre ses forces et ses besoins en matière d'apprentissage.

\textbf{nteragir avec le contenu} : Les parents peuvent explorer le contenu pédagogique proposé par \textbf{BEWIZE} et interagir avec les activités d'apprentissage de leur enfant. Ils peuvent consulter les leçons, les exercices et les ressources complémentaires pour avoir une idée précise de ce que leur enfant apprend. Cela leur permet de s'impliquer davantage dans le processus d'apprentissage de leur enfant et de lui fournir un soutien supplémentaire si nécessaire.

En offrant un accès facile à l'expérience d'apprentissage de leur enfant, \textbf{BEWIZE} permet aux parents d'être plus impliqués dans son éducation, de suivre ses avancées de près et de fournir un soutien personnalisé. Cela favorise une collaboration étroite entre les parents, les élèves et les enseignants pour assurer une expérience d'apprentissage enrichissante.

\subsection{Suivi en temps réel de la progression de votre enfant}
La progression de l'enfant est enregistrée en temps réel, offrant ainsi toutes les statistiques nécessaires, ainsi que des recommandations personnalisées en fonction de son niveau. Le suivi en temps réel permet d'accéder aux informations les plus récentes concernant les activités, les réponses et les progrès de l'enfant.

Des statistiques détaillées sont compilées pour donner un aperçu complet de la performance et de l'évolution de l'enfant. Les rapports clairs et complets permettent de comprendre ses forces, ses faiblesses, les domaines où il excelle et ceux qui nécessitent une attention particulière. Ces statistiques aident à mieux comprendre les progrès de l'enfant et à identifier les domaines où un soutien supplémentaire pourrait être nécessaire.

Des recommandations personnalisées sont fournies en se basant sur les données recueillies. Ces recommandations sont spécifiquement adaptées à l'enfant, et peuvent inclure des suggestions d'activités supplémentaires pour renforcer certaines compétences, des ressources complémentaires adaptées à son niveau, ou des conseils pédagogiques spécifiques pour l'aider à progresser de manière optimale.

Ce suivi en temps réel, les statistiques détaillées et les recommandations personnalisées offrent une vision complète de la progression de l'enfant. Cela permet aux parents de prendre des décisions éclairées pour soutenir et encourager leur enfant dans son parcours d'apprentissage.


%----------------------------------------------------------------------
%           Valeurs de \textbf{BEWIZE}
%----------------------------------------------------------------------
\section{Les Valeurs de BEWIZE}


En faisant cela, \textbf{BEWIZE s’appuie sur 7 valeurs principales} :

\subsection{Impact} 
BEWIZE s'engage à créer un impact pérenne dans le domaine de l'éducation en se concentrant sur l'amélioration de l'expérience d'apprentissage pour les générations futures. L'entreprise reconnaît que l'éducation est un pilier essentiel du développement individuel et collectif, et elle se donne pour mission de repenser et de réinventer les méthodes d'enseignement traditionnelles.

BEWIZE cherche à aller au-delà des approches éducatives conventionnelles en proposant des solutions innovantes et en adoptant une approche centrée sur l'apprenant. L'entreprise vise à fournir des expériences d'apprentissage qui sont plus engageantes, plus interactives et mieux adaptées aux besoins individuels des élèves. En utilisant les dernières avancées technologiques et en intégrant des méthodes pédagogiques éprouvées, BEWIZE cherche à transformer l'éducation et à créer des opportunités d'apprentissage significatives et durables.

En se concentrant sur l'impact à long terme, BEWIZE vise à préparer les apprenants pour l'avenir en développant les compétences et les connaissances nécessaires pour réussir dans un monde en constante évolution. L'entreprise reconnaît l'importance de développer des compétences telles que la pensée critique, la résolution de problèmes, la collaboration et la créativité, qui sont essentielles pour s'adapter aux défis et aux opportunités du 21e siècle.

De plus, BEWIZE cherche à combler les écarts dans l'accès à l'éducation en fournissant des solutions éducatives numériques qui sont facilement accessibles et disponibles pour tous. L'entreprise s'engage à rendre l'éducation plus inclusive en offrant des outils et des ressources qui répondent aux besoins des apprenants de tous les horizons.

En travaillant à créer un impact durable dans le domaine de l'éducation, BEWIZE aspire à laisser un héritage positif pour les générations futures. L'entreprise croit au pouvoir de l'éducation pour changer des vies, ouvrir des portes et contribuer au progrès de la société dans son ensemble. Par son engagement à repenser l'éducation et à créer des expériences d'apprentissage transformatrices, BEWIZE joue un rôle actif dans la construction d'un avenir meilleur pour tous.


\subsection{Partage} 
BEWIZE valorise le partage et la collaboration en tant que moteurs essentiels du progrès éducatif. L'entreprise s'engage à développer constamment en apprenant des autres et en partageant son savoir avec la communauté éducative.

BEWIZE reconnaît que l'éducation ne se limite pas à une seule entité, mais qu'elle est le fruit d'un échange constant entre enseignants, élèves, parents et autres acteurs du domaine éducatif. Dans cette optique, BEWIZE encourage activement la coopération entre les différents intervenants afin d'enrichir les connaissances et les pratiques pédagogiques.

L'entreprise favorise la création d'un écosystème éducatif dynamique en encourageant les enseignants à partager leurs idées, leurs expériences et leurs meilleures pratiques. BEWIZE organise des forums, des ateliers et des sessions de formation où les enseignants peuvent se rencontrer, échanger des idées et collaborer pour améliorer l'enseignement et l'apprentissage.

En partageant son savoir, BEWIZE contribue également à l'élargissement des connaissances de la communauté éducative. L'entreprise publie régulièrement du contenu éducatif de qualité, des articles, des guides et des ressources pédagogiques accessibles à tous. Elle encourage également les enseignants et les parents à partager leurs propres ressources et stratégies d'apprentissage, créant ainsi une culture de collaboration et d'entraide.

\subsection{Potentiel} 
BEWIZE croit fermement que comprendre l'autre est la clé pour libérer le potentiel de chaque individu. L'entreprise reconnaît que chaque élève a des talents uniques et des besoins spécifiques, et elle s'efforce de les comprendre et de les accompagner dans leur parcours éducatif.

BEWIZE s'inspire des autres pour accomplir sa mission et relever de grands défis. L'entreprise reconnaît que l'éducation ne se limite pas à une seule approche ou à une seule source de connaissances. Elle s'appuie sur les idées, les expériences et les perspectives des autres acteurs de l'éducation, tels que les enseignants, les chercheurs, les parents et les élèves eux-mêmes, pour enrichir ses propres pratiques et développer des solutions novatrices.

En comprenant les besoins et les aspirations de chaque élève, BEWIZE est en mesure de proposer une approche éducative personnalisée qui favorise leur épanouissement et leur réussite. L'entreprise met l'accent sur l'écoute active, la communication ouverte et la collaboration avec tous les acteurs impliqués dans l'éducation, afin de créer un environnement propice à l'apprentissage et au développement du potentiel de chaque individu.

En s'inspirant des autres et en reconnaissant l'importance de comprendre chaque élève, BEWIZE se positionne comme un partenaire qui encourage et soutient les élèves dans leur parcours éducatif. L'entreprise cherche à créer une synergie entre les différentes perspectives et compétences, afin de permettre à chaque élève de déployer pleinement son potentiel et de réaliser ses aspirations.

\subsection{Autonomie} 
BEWIZE met l'accent sur la capacité des individus à accomplir leur travail en toute liberté et à prendre des responsabilités. L'entreprise reconnaît que l'autonomie est un élément essentiel pour favoriser l'apprentissage efficace et le développement de compétences clés.

BEWIZE offre aux élèves un environnement d'apprentissage qui encourage l'autonomie. Grâce à son application mobile et à son back-office, les élèves ont accès à un contenu pédagogique interactif et peuvent progresser à leur propre rythme. Ils ont la liberté de choisir les sujets sur lesquels ils souhaitent se concentrer, d'explorer différents parcours d'apprentissage et de participer activement à leur propre éducation.

L'entreprise encourage également les élèves à prendre des responsabilités dans leur apprentissage. Ils sont incités à fixer des objectifs personnels, à suivre leur propre progression et à s'auto-évaluer. BEWIZE fournit des outils et des ressources qui permettent aux élèves de mesurer leurs propres compétences, de réfléchir sur leurs réalisations et de prendre des décisions éclairées sur les domaines dans lesquels ils souhaitent se développer davantage.

En mettant l'autonomie au cœur de son approche, BEWIZE vise à développer chez les élèves des compétences essentielles telles que la prise d'initiative, la gestion du temps, la prise de décision et la responsabilité personnelle. L'entreprise croit en la capacité des élèves à être acteurs de leur propre apprentissage, et elle fournit les outils et les ressources nécessaires pour les soutenir dans cette démarche.

En encourageant l'autonomie, BEWIZE permet aux élèves de développer leur confiance en eux, leur motivation intrinsèque et leur capacité à s'adapter aux défis futurs. Ils acquièrent également des compétences transversales telles que la gestion de l'information, la résolution de problèmes et la pensée critique, qui sont essentielles pour réussir dans un monde en constante évolution.

\subsection{Épanouissement} 
BEWIZE est animé par une passion profonde pour l'apprentissage, la recherche et l'innovation. L'entreprise reconnaît que l'épanouissement personnel et professionnel est étroitement lié à la volonté de s'enrichir continuellement et de partager ses connaissances avec les autres.

BEWIZE encourage ses membres à embrasser une culture d'apprentissage permanent. Les équipes sont encouragées à se tenir au courant des dernières avancées dans le domaine de l'éducation, à suivre des formations et des programmes de développement professionnel, et à s'engager dans des projets de recherche et d'innovation. L'entreprise encourage également ses membres à participer à des conférences, à des ateliers et à d'autres événements éducatifs pour élargir leurs horizons et enrichir leur compréhension des meilleures pratiques.

En se consacrant à l'enrichissement de son propre savoir, BEWIZE est en mesure de proposer des solutions éducatives innovantes et à la pointe de la technologie. L'entreprise investit dans la recherche et le développement pour repousser les limites de l'apprentissage numérique et créer des expériences d'apprentissage engageantes, interactives et personnalisées.

BEWIZE croit également en l'importance de partager ses connaissances avec les autres. Les membres de l'entreprise sont encouragés à partager leurs découvertes, leurs idées et leurs bonnes pratiques avec la communauté éducative. Ils peuvent contribuer à la rédaction d'articles, à la création de ressources pédagogiques, à la participation à des groupes de discussion et à d'autres activités de partage des connaissances. Cette volonté de partager ses connaissances favorise la croissance collective et contribue à l'enrichissement de l'éducation dans son ensemble.

En s'épanouissant par l'apprentissage, la recherche et l'innovation, BEWIZE se positionne comme un acteur engagé et passionné dans le domaine de l'éducation. L'entreprise cherche constamment à évoluer, à repousser les limites et à transformer la façon dont nous apprenons, pour le bénéfice de tous les apprenants et de la société dans son ensemble.

\subsection{Innovation} 
BEWIZE se positionne comme un leader dans le domaine de l'innovation éducative. L'entreprise est constamment à la recherche de moyens créatifs pour éliminer les obstacles, résoudre les problèmes et générer des solutions pérennes qui répondent aux besoins changeants du domaine de l'éducation.

BEWIZE croit fermement que l'innovation est essentielle pour transformer l'expérience d'apprentissage et améliorer les résultats des élèves. L'entreprise encourage ses équipes à repousser les limites de la pensée traditionnelle et à adopter une approche audacieuse et novatrice. Elle encourage la recherche de nouvelles idées, de nouvelles méthodes et de nouvelles technologies qui peuvent révolutionner l'éducation.

Les équipes de BEWIZE sont encouragées à collaborer avec des experts externes, à nouer des partenariats avec des institutions académiques et à rester à l'affût des dernières tendances et avancées dans le domaine de la technologie éducative. L'entreprise investit également dans la recherche et le développement pour créer des solutions éducatives innovantes qui sont adaptées aux besoins spécifiques des apprenants d'aujourd'hui.

En cherchant continuellement à innover, BEWIZE s'efforce d'identifier et de résoudre les obstacles auxquels sont confrontés les élèves et les éducateurs. Que ce soit en repensant les méthodes d'enseignement, en développant de nouvelles plateformes numériques ou en intégrant des technologies émergentes telles que l'intelligence artificielle et la réalité virtuelle, l'entreprise s'engage à offrir des solutions qui ont un impact réel et durable.

En mettant l'innovation au cœur de ses activités, BEWIZE aspire à transformer l'éducation en créant des expériences d'apprentissage engageantes, interactives et adaptées aux besoins individuels des élèves. L'entreprise cherche à repousser les limites de ce qui est possible dans l'éducation et à ouvrir de nouvelles perspectives pour les apprenants, les éducateurs et l'ensemble du système éducatif.

\subsection{Passion} 
La passion est au cœur de toutes les activités de BEWIZE. Elle est la force motrice qui inspire l'excellence et pousse l'entreprise à viser la perfection dans tout ce qu'elle entreprend. Chez BEWIZE, la passion est présente à tous les niveaux, des équipes de développement et de conception jusqu'aux équipes de support et de gestion.

La passion de BEWIZE se reflète dans la qualité exceptionnelle de ses solutions éducatives. L'entreprise s'efforce de créer des expériences d'apprentissage qui sont non seulement efficaces sur le plan pédagogique, mais aussi captivantes et engageantes pour les apprenants. Chaque aspect des produits et services de BEWIZE est conçu avec soin et attention aux détails, dans le but de fournir des solutions éducatives de la plus haute qualité.

La passion de BEWIZE se traduit également par son engagement envers l'innovation continue. L'entreprise ne se contente pas de suivre les tendances, mais cherche à être à l'avant-garde de l'évolution de l'éducation. Elle explore de nouvelles idées, teste de nouvelles approches et utilise les dernières avancées technologiques pour offrir des solutions éducatives innovantes et efficaces.

Cette passion se reflète également dans l'engagement de BEWIZE envers la réussite des apprenants. L'entreprise est profondément investie dans le progrès et le développement des élèves, et elle s'efforce de fournir les outils et les ressources nécessaires pour les aider à atteindre leur plein potentiel.

En mettant en œuvre une approche passionnée et dévouée, BEWIZE se distingue en tant qu'acteur majeur dans le domaine de l'éducation. L'entreprise incarne une culture de l'excellence et de l'engagement envers l'amélioration continue, plaçant la passion au cœur de son travail pour créer des solutions éducatives qui inspirent et transforment la vie des apprenants.








%----------------------------------------------------------------------
%           Activités principales de l'entreprise
%----------------------------------------------------------------------
\section{Activités principales de l'entreprise}
\subsection{Développement de l'application mobile Bewize}  
    
    Le développement de l'application mobile Bewize est réalisé par une équipe de professionnels hautement qualifiés en programmation, en design d'interface utilisateur et en développement d'applications mobiles. Ces experts mettent en œuvre leurs compétences techniques pour créer une application robuste et conviviale, qui offre une expérience d'apprentissage engageante pour les élèves.

    L'équipe de développement de Bewize travaille en étroite collaboration avec les experts pédagogiques de l'entreprise pour concevoir une interface intuitive et interactive. L'objectif est de rendre l'application accessible aux élèves de tous les niveaux scolaires et de tous les profils d'apprentissage. Ils veillent à ce que l'application offre une expérience fluide et agréable, tout en favorisant l'autonomie et l'épanouissement des élèves.
    
    L'application mobile Bewize propose une variété de fonctionnalités adaptées aux besoins pédagogiques. Les élèves peuvent accéder à des quiz interactifs qui leur permettent de tester leurs connaissances et de consolider leurs apprentissages. Des jeux éducatifs ludiques sont également intégrés pour rendre l'apprentissage plus amusant et stimulant. De plus, l'application offre un accès facile aux ressources pédagogiques, telles que des leçons, des exercices et des supports de cours, qui sont alignés sur le référentiel du ministère de l'Éducation.
    
    L'équipe de développement travaille en permanence pour améliorer et mettre à jour l'application mobile Bewize. Ils prennent en compte les retours des utilisateurs et s'adaptent aux dernières avancées technologiques pour garantir une expérience optimale. La sécurité des données des utilisateurs est une priorité absolue, et des mesures de protection sont mises en place pour assurer la confidentialité des informations personnelles.
    
    En combinant expertise technique et connaissances pédagogiques, l'équipe de développement de Bewize crée une application mobile qui offre un environnement d'apprentissage interactif et personnalisé pour les élèves. Grâce à cette application, les élèves peuvent explorer les contenus pédagogiques, suivre leur progression en temps réel et bénéficier de recommandations personnalisées, contribuant ainsi à leur réussite éducative.
    
\subsection{Création du tableau de bord école (back-office)} 
    
    La création du tableau de bord école (back-office) est une autre composante clé du travail de Bewize. Ce tableau de bord est spécifiquement conçu pour les enseignants, leur fournissant un accès centralisé à toutes les informations nécessaires pour suivre et gérer les progrès des élèves.

    L'équipe de professionnels de Bewize dédiée à ce projet met en œuvre son expertise en conception et développement de logiciels pour créer un tableau de bord école convivial et intuitif. Ils utilisent des technologies adaptées aux besoins spécifiques des enseignants, en veillant à ce que l'interface soit facile à naviguer et à comprendre.
    
    Le tableau de bord école permet aux enseignants de suivre les progrès des élèves en temps réel. Ils peuvent consulter les résultats des tests et des évaluations, analyser les performances individuelles et collectives, et identifier les domaines où les élèves peuvent avoir besoin d'un soutien supplémentaire. Ces informations leur permettent de fournir un suivi personnalisé, en adaptant leur enseignement aux besoins et aux capacités de chaque élève.
    
    Le back-office offre également des fonctionnalités avancées, telles que la génération de rapports et de statistiques détaillées. Les enseignants peuvent obtenir une vue d'ensemble des performances de la classe, identifier les tendances et les points forts, et ajuster leur approche pédagogique en conséquence. Cela leur permet de prendre des décisions éclairées pour améliorer l'apprentissage des élèves.
    
    L'équipe de développement travaille en étroite collaboration avec les enseignants et les experts pédagogiques de Bewize pour s'assurer que le tableau de bord école répond aux besoins réels des enseignants. Ils prennent en compte les retours d'utilisateurs et effectuent des mises à jour régulières pour améliorer les fonctionnalités et l'expérience globale.
    
    Grâce à ce tableau de bord école, Bewize offre aux enseignants un outil puissant pour gérer et optimiser l'apprentissage des élèves. Il leur permet de bénéficier d'un suivi personnalisé, de prendre des décisions basées sur des données et de collaborer de manière plus efficace avec les autres acteurs de l'éducation.
    
\subsection{Gestion de contenu}
    
    Bewize propose également un autre back-office dédié à la gestion de contenu au sein de l'application Bewize. Cette tâche consiste à créer, organiser et mettre à jour le contenu éducatif proposé aux élèves. Les professionnels de Bewize travaillent avec des experts pédagogiques pour développer du contenu de qualité, adapté aux programmes scolaires et aux besoins des élèves. Ils veillent à ce que le contenu soit à jour, pertinent et engageant.
    
\subsection{Suivi et support technique}
    
   Bewize attache une grande importance à offrir un suivi et un support technique de qualité à ses utilisateurs. Les professionnels de Bewize sont disponibles pour répondre aux questions, résoudre les problèmes techniques et fournir une assistance personnalisée pour garantir une expérience utilisateur optimale.

    L'équipe de support technique de Bewize est composée de professionnels qualifiés et expérimentés dans le domaine des technologies de l'information et de la communication. Ils sont formés pour traiter efficacement les demandes des utilisateurs et fournir des solutions adaptées à leurs besoins.
    
    Lorsque les utilisateurs rencontrent des problèmes techniques, tels que des difficultés de connexion, des erreurs d'affichage ou des questions sur les fonctionnalités de l'application Bewize, l'équipe de support technique est disponible pour les aider. Ils peuvent guider les utilisateurs à travers les étapes de résolution des problèmes, fournir des explications claires et résoudre les problèmes de manière rapide et efficace.
    
    En plus de répondre aux questions et de résoudre les problèmes, l'équipe de support technique de Bewize est également chargée de recueillir les commentaires des utilisateurs. Ils sont à l'écoute des retours d'expérience des utilisateurs, de leurs suggestions d'amélioration et de leurs besoins spécifiques. Ces informations sont ensuite utilisées pour effectuer des mises à jour et des améliorations de l'application Bewize, garantissant ainsi une évolution continue de la plateforme en fonction des besoins réels des utilisateurs.
    
    Grâce à un suivi et un support technique attentifs, Bewize s'assure que les utilisateurs de l'application bénéficient d'une expérience fluide, sans interruption majeure et répondant à leurs attentes. L'équipe de support technique joue un rôle clé dans la satisfaction des utilisateurs et dans l'optimisation de l'utilisation de l'application Bewize dans les environnements éducatifs.
    
\subsection{Collaboration avec les établissements scolaires}
    
    La collaboration avec les établissements scolaires est une étape essentielle dans le déploiement de l'application Bewize. Bewize reconnaît l'importance de travailler en étroite collaboration avec les enseignants et les responsables d'établissement pour intégrer efficacement l'application dans les salles de classe.

    Dans le cadre de cette collaboration, l'équipe de Bewize prend contact avec les établissements scolaires afin de présenter l'application Bewize et d'expliquer ses fonctionnalités et ses avantages pédagogiques. Ils mettent en évidence comment l'application peut améliorer l'expérience d'apprentissage des élèves et soutenir les enseignants dans leur pratique pédagogique.
    
    Une fois l'intérêt suscité, l'équipe de Bewize propose des formations aux enseignants et au personnel de l'établissement. Ces formations sont conçues pour familiariser les utilisateurs avec l'application, ses fonctionnalités et ses possibilités d'utilisation en classe. L'objectif est d'assurer une utilisation optimale de l'application Bewize, en tenant compte des besoins et des objectifs pédagogiques spécifiques de chaque établissement.
    
    Pendant la phase de déploiement, l'équipe de Bewize fournit un soutien continu aux établissements scolaires. Ils sont disponibles pour répondre aux questions, fournir des conseils et résoudre les éventuels problèmes rencontrés lors de l'utilisation de l'application. Ce soutien permet aux enseignants de bénéficier d'une assistance personnalisée tout au long de leur parcours d'intégration de l'application Bewize dans leurs pratiques d'enseignement.
    
    La collaboration avec les établissements scolaires permet à Bewize de mieux comprendre les besoins et les attentes des enseignants et des élèves. Cela permet également de recueillir des commentaires et des retours d'expérience précieux, qui sont utilisés pour continuer à améliorer l'application et adapter les fonctionnalités en fonction des besoins réels des utilisateurs.
    
    En travaillant en étroite collaboration avec les établissements scolaires, Bewize assure une intégration harmonieuse de l'application Bewize dans les pratiques éducatives, maximisant ainsi son impact et sa valeur pédagogique pour les élèves et les enseignants.

%----------------------------------------------------------------------
%              Fiche signalétique de Bewize
%----------------------------------------------------------------------
\section{Fiche signalétique de Bewize}
\begin{center}
    
\begin{tabular}{ |p{7.5cm}|p{3cm}| }
 
 \hline
 Raison sociale & Bewize \\
\hline
 Siège social & — \\
 \hline
 Forme juridique & SARL AU \\
 \hline
 Contact & Mohamed Mekran \\
 \hline
 Date de création & Août 2022 \\
 \hline
 Capital social & 100K DHS \\
 \hline
 Effectif & 6 \\
 \hline
 Secteur d’activité & Edtech \\
 \hline

\end{tabular}
\end{center}



%----------------------------------------------------------------------
%               Structure de Bewize    
%----------------------------------------------------------------------
\section{Structure de Bewize}
\begin{figure}[H]
    \includegraphics[width=1\textwidth]{Figures/structure_de_bewize.png}\\[0.5cm]
    \centering
    \caption{Structure de BEWIZE}
\end{figure}

\begin{itemize}
    \item[$\bullet$] \textbf{CEO (Directeur Général):}  est le dirigeant de l'entreprise Bewize. Il est responsable de la gestion globale de l'entreprise, de la définition de la vision et des objectifs stratégiques, ainsi que de la prise de décisions importantes. Il supervise également les différentes équipes et assure la coordination entre les différents départements. \\
    
    \item[$\bullet$] \textbf{Product Manager (Chef de produit):} Le Product Manager est responsable de la gestion et du développement des produits de Bewize. Il travaille en étroite collaboration avec les équipes techniques et marketing pour définir les fonctionnalités et les caractéristiques des produits, réaliser des études de marché, suivre les tendances du secteur et s'assurer que les produits répondent aux besoins des clients.\\
    
    \item[$\bullet$] \textbf{Digital Marketer (Responsable marketing numérique):}
    Le Digital Marketer est chargé de la promotion et de la commercialisation des produits de Bewize sur les canaux en ligne. Il met en œuvre des stratégies de marketing numérique, telles que la publicité en ligne, le référencement, les médias sociaux, le marketing par e-mail, afin d'attirer des prospects et de convertir les clients. Le Digital Marketer analyse également les performances des campagnes marketing et propose des améliorations.\\

    \item[$\bullet$] \textbf{Graphic Designer (Graphiste) :}  Le Graphic Designerest responsable de la création visuelle et graphique des produits et des supports de communication de Bewize. Il conçoit des éléments graphiques attrayants tels que des logos, des illustrations, des interfaces utilisateur, des bannières publicitaires, des infographies, etc. Le Graphic Designer travaille en étroite collaboration avec l'équipe marketing pour créer une identité visuelle cohérente et attrayante pour l'entreprise.\\
    
    \item[$\bullet$] \textbf{ Software Engineer (Ingénieur logiciel) :}  Le Software Engineer est responsable du développement et de la maintenance des logiciels et des systèmes informatiques de Bewize. Il programme et code les fonctionnalités des produits, effectue des tests, résout les problèmes techniques et optimise les performances. Le Software Engineer collabore également avec l'équipe Product Manager pour comprendre les besoins des utilisateurs et mettre en œuvre les fonctionnalités requises.\\
    
    \item[$\bullet$] \textbf{ Responsable Pédagogique :} Le Responsable Pédagogique est chargé de superviser les aspects liés à la pédagogie au sein de Bewize. Il est responsable de la conception des programmes de formation, de l'élaboration des contenus pédagogiques, du recrutement des formateurs et de l'évaluation de l'efficacité des programmes. Le Responsable Pédagogique s'assure que les programmes de formation de Bewize répond
    aux besoins du marché et des apprenants, tout en garantissant la qualité des cours dispensés. \\

    Cette structure organisationnelle reflète une approche multidisciplinaire de Bewize, où chaque membre de l'équipe apporte des compétences spécifiques pour contribuer au succès de l'entreprise. La collaboration entre ces différents rôles permet de développer, promouvoir et améliorer les produits et services de Bewize, tout en assurant une expérience client satisfaisante.


\end{itemize}
